%% The seven-class (v2) text for mars_global_mosaic.tex, drafted before any v2 number exists (NEXT-STEPS §12.4 W8,
%% KB §69.18). NOT \input by the paper: at the freeze, copy the branch that D18's verdict picks into the place named
%% above each block, replace every \VV{log}{key} with the number from that log (build\logs\, written by run_v2.ps1;
%% percentages to one decimal, kappa to three), put each number into the KB first, then run verify_interim.py on the
%% .tex. Drafted against the rehearsal's logs (*_v2p*), whose numbers are never results and appear nowhere below.
%%
%% The logs, by the queue's step names (run_v2.ps1):
%%   SV2   03_score_svm_v2           global60_classification_check_v2_200m_mode11.json
%%   SF2   04_score_svm_four_on_v2   global60_classification_check_200m_mode11_on_v2_labels.json
%%   D18   05_d18_verdict_v2         d18_verdict.json      (replaces: true/false picks the branch)
%%   GEO   06_geomap_v2              geomap_check_v2.json
%%   UV2   12_score_unet_v2          global60_classification_check_global60_landforms_unet_v2_200m_seamfix.json
%%   UF2   13_score_unet_four_on_v2  global60_classification_check_global60_landforms_unet_200m_seamfix_on_v2_labels.json
%%   DU    not in the queue: score_d18_verdict.py --four <UF2> --v2 <UV2> --out logs\d18_verdict_unet_v2.json
%%         (the rehearsal's d18_verdict_unet_v2p.json was made the same way)
%%   MET   score_metrics.py logs\<SV2> logs\<SF2> logs\<UV2> logs\<UF2> --out logs\score_metrics_v2.json
%% Keys are as the scorers write them: held_out.accuracy, held_out.producers.<class>, d18 a_accuracy_and_kappa.
%% accuracy = [v2, four], c_kept_classes.<class>.fall, and so on.
%%
%% Class colours, the same in layouts, deck and paper: labelsets.COLOURS (the labels' own; checked 2026-10-27 under
%% simulated protan, deutan and tritan vision, closest pair Delta E 21.1, KB §69.18). Do not pick new ones.
\newcommand{\VV}[2]{\textbf{[#1: #2]}}   % only for a trial compile of this file's text

%% ---------------------------------------------------------------------------------------------------------
%% 1. Methods, "Supervised classification and its score": a paragraph after the four-class one. Either branch.
%% ---------------------------------------------------------------------------------------------------------
An audit of the 512 labels then revised the scheme to seven classes: Crater narrowed to impact craters, volcanic
plains drawn as Normal Ground, and three new classes, volcanic edifice, apron and fluvial channel. Proposals for each
label and new polygons for the new classes were built from the geologic map, the Robbins database, the IAU names, the
elevation model and the valley map of \citet{Alemanno2018}. Every relabel, check and edifice outline was reviewed by
eye; the apron and fluvial-strip proposals were screened by a rule written before it ran. The result, 971 labels, keeps
the same 14 held-out blocks (734 training and 237 held-out polygons), and a support vector machine and a U-Net were
retrained on it exactly as before. The seven-class map was to replace the published one only if, on the same v2
held-out polygons, (a) its accuracy and \kap\ on the ground both maps can name (v2 classes 1--4) were at least the
published map's, (b) each new class had a user's accuracy above its share of the held-out area, and (c) Crater, hills
and Normal Ground each lost under 5 points of producer's accuracy. %% KB §68.1--§68.4, §69.2, §69.8

%% ---------------------------------------------------------------------------------------------------------
%% 2a. Results, a new subsection "Seven classes" after "The landform classification". BRANCH: D18 replaces = true.
%% ---------------------------------------------------------------------------------------------------------
On the 237 v2 held-out polygons the seven-class support vector machine map scores \VV{SV2}{held\_out.accuracy},
\kap\,\VV{SV2}{held\_out.kappa}, against \VV{SF2}{held\_out.accuracy}, \kap\,\VV{SF2}{held\_out.kappa} for the published
map, which cannot name the new classes. On the ground both can name it scores \VV{D18}{a.accuracy[0]} against
\VV{D18}{a.accuracy[1]}; volcanic edifice, apron and fluvial channel reach user's accuracies of
\VV{D18}{b.volcanic edifice.users}, \VV{D18}{b.apron.users} and \VV{D18}{b.fluvial channel.users} against area shares
of \VV{D18}{b.*.area\_share}; and no kept class loses 5 points (largest fall \VV{D18}{c.max fall}, \VV{D18}{c.class}).
By the rule written before the run the seven-class map replaces the four-class one as the published landform map
(Figure~\ref{fig:landforms7}). Its balanced accuracy is \VV{MET}{SV2 balanced\_accuracy} and its disagreement
\VV{MET}{SV2 quantity} quantity and \VV{MET}{SV2 allocation} allocation. Against the geologic map, volcanic edifice is
\VV{GEO}{enrichment edifice on volcanic edifice units} times as common on the volcanic edifice units as across \pmsixty.

%% ---------------------------------------------------------------------------------------------------------
%% 2b. The same place. BRANCH: D18 replaces = false. Name every failed part, as the verdict prints it.
%% ---------------------------------------------------------------------------------------------------------
On the 237 v2 held-out polygons the seven-class support vector machine map scores \VV{SV2}{held\_out.accuracy},
\kap\,\VV{SV2}{held\_out.kappa}, against \VV{SF2}{held\_out.accuracy}, \kap\,\VV{SF2}{held\_out.kappa} for the published
map. By the rule written before the run it does not replace it: \VV{D18}{the failed parts, each with its numbers, e.g.
"(a) on classes 1--4 it scores x against y"; "(b) apron's user's accuracy x is below its area share y"; "(c) Crater
loses x points"}. The four-class map stays the published landform map, and the seven-class map is reported as a test
(Figure~\ref{fig:landforms7}). What it shows is \VV{SV2}{one sentence on where the new classes land: producer's of
edifice, apron and fluvial channel}.

%% ---------------------------------------------------------------------------------------------------------
%% 3. Results, "Deep learning": after the U-Net paragraph. Either branch; DU decides the wording of its last sentence.
%% ---------------------------------------------------------------------------------------------------------
Retrained on the seven classes, the U-Net scores \VV{UV2}{held\_out.accuracy}, \kap\,\VV{UV2}{held\_out.kappa} on the
v2 held-out polygons, against \VV{UF2}{held\_out.accuracy} for the four-class U-Net, with new-class user's accuracies of
\VV{UV2}{users edifice / apron / fluvial channel}. %% DU replaces = true:
By the same rule it replaces the four-class U-Net as the province map.
%% DU replaces = false:
%% By the same rule it does not replace the four-class U-Net: \VV{DU}{the failed parts}.

%% ---------------------------------------------------------------------------------------------------------
%% 4. Table 3 (tab:svm): rows below the U-Net. The four-class rows keep their 183-polygon scores; say in the note
%%    that these rows are on the 237 v2 polygons and are not comparable with the rows above.
%% ---------------------------------------------------------------------------------------------------------
%% \hline
%% \multicolumn{7}{l}{\textit{On the 237 v2 held-out polygons}} \\
%% SVM four-class (published) & \VV{SF2}{acc} & \VV{SF2}{kappa} & \VV{MET}{bal} & \VV{MET}{F1} & \VV{MET}{Q} & \VV{MET}{A} \\
%% SVM seven-class & \VV{SV2}{acc} & \VV{SV2}{kappa} & \VV{MET}{bal} & \VV{MET}{F1} & \VV{MET}{Q} & \VV{MET}{A} \\
%% U-Net four-class & \VV{UF2}{acc} & ... \\
%% U-Net seven-class & \VV{UV2}{acc} & ... \\

%% ---------------------------------------------------------------------------------------------------------
%% 5. Figure: the seven-class map, from the new layout beside 04 (NEXT-STEPS §11.5 step 4), and the per-class chart
%%    svm_vs_unet.png redrawn on v2 (make_fig_final.py). Label: fig:landforms7.
%% ---------------------------------------------------------------------------------------------------------

%% ---------------------------------------------------------------------------------------------------------
%% 6. Discussion, "Limits": replace the sentence "Both maps are scored on the same held-out polygons, and the
%%    seven-class map replaces the published one only by a rule written before the run." with one branch, and
%%    replace the held-out paragraph's first sentence.
%% ---------------------------------------------------------------------------------------------------------
%% replaces = true:
The seven-class map passed the rule set before it ran and is the published map; its apron and fluvial labels rest on
a written screen, not on inspection, and its hills class is short of the training minimum.
%% replaces = false:
%% The seven-class map failed the rule set before it ran (\VV{D18}{the failed part}), so the four-class map, with its
%% volcano-shaped Crater labels, remains the published one.
%% held-out use:
The held-out blocks have now been used twice, once by the four-class models and once by the seven-class ones; any
further tuning needs validation blocks of its own.

%% ---------------------------------------------------------------------------------------------------------
%% 7. Abstract and Conclusions: one clause each, in the 250-word budget (the abstract is at 250 now: cut as much).
%% ---------------------------------------------------------------------------------------------------------
%% replaces = true:  "Revised to seven classes after an audit of the labels, the map scores \VV{SV2}{acc} on 237
%%                    held-out polygons and replaces the four-class map by a rule set in advance."
%% replaces = false: "A seven-class revision after an audit of the labels scores \VV{SV2}{acc} on 237 held-out
%%                    polygons and fails the rule set in advance to replace the four-class map."
